\section{AC (Small-Signal) Analysis}
\label{sec:ac-analysis}

The solver of Section~\ref{sec:circuit-solver} answers a transient question: given a
circuit's exact values and a starting state, what does every voltage and current do over
time? A complementary question -- how does a circuit respond to a sinusoidal excitation as a
function of frequency alone -- is answered instead by AC (small-signal) analysis, the
technique underlying every Bode plot in an electronics textbook: every node voltage and
branch current is assumed already varying sinusoidally at a single angular frequency
$\omega$, and the solver solves directly for each phasor's complex amplitude, using the same
node/branch structure as the transient MNA formulation but with every conductance,
companion model, and source replaced by its frequency-domain counterpart. CircuitCraft
implements this as an entirely separate solver, \texttt{AcCircuit}, sharing the transient
solver's node-numbering topology (Section~\ref{sec:architecture}) but operating on disjoint,
complex-valued representations solved fresh at each swept frequency. Following the standard
small-signal convention, every independent voltage source other than the one designated as
the AC excitation is stamped at exactly zero volts (an ideal short) for the sweep's duration,
isolating the frequency response to the one signal path being characterized.

Each element gains a second stamping method taking $\omega$ directly: a resistor's admittance
is its unchanged conductance $1/R$; a capacitor's is $Y=j\omega C$ and an inductor's
$Y=1/(j\omega L)$, requiring no companion history term since AC analysis solves the
steady-state response directly; a diode linearizes about its last-solved operating point into
the small-signal resistance $r_d=V_T/I_D$; and a memristor is, for now, treated as a plain
resistor frozen at whatever resistance it currently holds, a deliberate simplification
(Section~\ref{sec:limitations}) rather than an oversight -- a genuine small-signal memristor
model would need its own frequency response derived from the charge-controlled state
equation, and recent results toward exactly such a model, including Bode-plot-based
extraction from real flexible-memristor measurements, are given
in~\cite{gater2026smallsignal,gater2025bodeplots}. Each transient-verified case above was
independently checked against a closed-form AC reference during development
(Section~\ref{sec:validation}): a resistive divider's response is flat across many decades of
frequency, and a series RC low-pass reproduces the textbook $-3$~dB, $-45^\circ$ point at
$f_c=1/(2\pi RC)$ with the correct asymptotic $-20$~dB/decade rolloff.

A new two-terminal AC Source component, wired like any other, supplies the sweep's
excitation; unlike every other adjustable component, its amplitude and swept frequency range
are set only through the shift-right-click value editor rather than a preset cycle, since a
frequency \emph{range} is not the kind of quantity a short preset list represents well. In
the ordinary transient simulation it contributes nothing but a 0~V source, electrically a
plain wire, so a circuit built around one behaves normally when no sweep is running. The
ideal op-amp, an infinite-bandwidth nullor in the transient model, is replaced during a sweep
by a standard two-pole open-loop gain model with fixed, illustrative values ($A_0=10^5$,
poles at 20~Hz and 3~MHz) giving it the rolloff and phase shift a real amplifier -- and a
Bode plot -- would show, verified directly against the model's own closed-form $-3$~dB point.

A dedicated AC probe captures a two-step interaction distinct from every other probe:
right-clicking an AC Source first pins it as the sweep's excitation; right-clicking a second
point afterward -- any wired component, wire, or ground block -- runs the sweep and displays
the result. The sweep evaluates 60 log-spaced frequencies across the pinned source's
configured range, solving a fresh \texttt{AcCircuit} at each one and reporting the magnitude
(dB) and phase (degrees) of the signal position's voltage relative to the source's own
amplitude, rendered as two stacked graphs against a shared log-frequency axis.
